\section{The two-transitive envelope}\label{sec:action}

We now turn the algebraic construction into the required permutation
representation.  Proposition~\ref{prop:affine} and
Lemma~\ref{lem:transition} give the finitely presented group
\[
 \Gamma_0=B^6\rtimes\St_6(R)
\]
and its endomorphism $\Theta=(\phi^{(6)},\theta)$.
Theorem~\ref{thm:kernel} says that $\theta$ kills the kernel of the matrix
map.  A mapping torus will remove this kernel while retaining finite
presentation.  We then use the simple subalgebra $A\subseteq B$ to obtain
two-transitivity.

\subsection{The mapping torus and its linear part}

We first record explicitly why a mapping torus has the needed description
even when its defining endomorphism is not injective.

\begin{lemma}\label{lem:mapping-torus}
Let $Q$ be a group and $f:Q\to Q$ an endomorphism.  Write
$L=\varinjlim(Q,f)$, with classes $[q,j]$ for $j\geq0$ and the convention
$[q,j]=[f(q),j+1]$.  The map
\[
 \alpha([q,j])=[q,j+1]
\]
is an automorphism of $L$, and
\[
 \langle Q,t\mid t^{-1}qt=f(q)\ (q\in Q)\rangle
 \cong L\rtimes_\alpha\mathbb Z.
\]
If $Q$ is finitely presented, so is this group.
\end{lemma}

\begin{proof}
Equality in the direct limit means equality after sufficiently many
forward transitions.  Thus $\alpha$ is well defined, as is
$\beta([q,j])=[f(q),j]$.  The defining equivalence gives
$\alpha\beta=\beta\alpha=\mathrm{id}$.
Let $t$ act on $L$ by $\alpha$.  Then
\[
 t^{-1}[q,0]t=\beta([q,0])=[f(q),0],
\]
so the displayed presentation maps to $L\rtimes_\alpha\mathbb Z$.
In the other direction, send
\[
 [q,j]\longmapsto t^jqt^{-j}.
\]
The relation $t^{-1}qt=f(q)$ implies
$t^{j+1}f(q)t^{-(j+1)}=t^jqt^{-j}$, proving that this map is well
defined.  It intertwines $\alpha$ with conjugation by $t$, and hence
extends to a homomorphism from the semidirect product.  The two maps are
inverse on $Q$ and $t$.

For a finite presentation $Q=\langle X\mid\mathcal R\rangle$, choose a
word representing $f(x)$ for each $x\in X$.  Adjoin $t$ and the finitely
many relations $t^{-1}xt=f(x)$ for $x\in X$.  These imply the required
relations for every element of $Q$, since both conjugation and $f$ are
homomorphisms.
\end{proof}

Apply Lemma~\ref{lem:mapping-torus} to $(\Gamma_0,\Theta)$ and denote
the resulting finitely presented group by $\Gamma$.  Put
\[
 M_\infty=\varinjlim((B,+),\phi),\qquad
 E_\infty=\varinjlim(\E_6(R),F),
 \qquad F(g)=I+\psi^{(6)}(g-I).
\]
The transition $F$ is the injective homomorphism of
Lemma~\ref{lem:transition}; recall the matrix map
$\pi:\St_6(R)\to\E_6(R)$.

\begin{lemma}\label{lem:affine-limit}
There are compatible isomorphisms
\[
 \varinjlim(\Gamma_0,\Theta)
 \cong M_\infty^6\rtimes E_\infty,
 \qquad
 \Gamma\cong(M_\infty^6\rtimes E_\infty)\rtimes\mathbb Z.
\]
In the second isomorphism, conjugation by $t$ shifts the stage index
forward by one on both $M_\infty$ and $E_\infty$.
\end{lemma}

\begin{proof}
The square $\pi\theta=F\pi$ induces a surjection from the
Steinberg direct limit to $E_\infty$.  If a class represented by
$h\in\St_6(R)$ maps to the identity, some iterate of $F$ sends
$\pi(h)$ to $I$.  Since $F$ is injective, $\pi(h)=I$ already.
Theorem~\ref{thm:kernel} now gives $\theta(h)=1$, so that class is
trivial.

The compatibility of the module action is
\begin{equation}\label{eq:action-transition}
 \phi^{(6)}(g v)=F(g)\phi^{(6)}(v)
 \qquad(g\in\E_6(R),\ v\in B^6).
\end{equation}
Indeed $\phi(r\cdot b)=\psi(r)\cdot\phi(b)$ for $r\in R$ and
$b\in B$, and separating $g$ into $I+(g-I)$ gives exactly the
modified-identity formula for $F$.  Thus matrices and vectors act on one
another after passage to any common stage.

The stage maps
\[
 (v,h)\longmapsto([v,j],[\pi(h),j])
\]
therefore induce a homomorphism from $\varinjlim(\Gamma_0,\Theta)$ to
$M_\infty^6\rtimes E_\infty$.  It is surjective: choose a common stage
for a vector and a matrix, and then choose a Steinberg preimage of that
matrix.  If a representative $(v,h)$ maps to the identity, injectivity
of $\phi$ gives $v=0$, and injectivity of $F$ gives $\pi(h)=I$.
Consequently $\Theta(v,h)=(0,1)$ by Theorem~\ref{thm:kernel}.  The
homomorphism is injective.  The construction respects stage shifts, so
Lemma~\ref{lem:mapping-torus} gives the second isomorphism as well.
\end{proof}

\subsection{A faithful action and two-transitivity}

Set $X=M_\infty^6$. The subgroup $M_\infty^6$ acts on $X$ by
translations, and $E_\infty$ acts by matrices after passing to a common
stage, as in \eqref{eq:action-transition}. Let $t$ act by the forward
stage shift. This shift intertwines the matrix actions, so these formulas define an
action of $\Gamma$ on $X$.

\begin{proposition}\label{prop:two-transitive-action}
The set $X$ is countably infinite.  The action of $\Gamma$ on $X$ is
faithful and two-transitive, and all its point stabilizers are finitely
generated.
\end{proposition}

\begin{proof}
\emph{The stage filtration.}
Since $\phi$ is injective, identify the $j$th additive stage with its
image $M_j\subseteq M_\infty$, and put $X_j=M_j^6$.  The stages are
nested because $[b,j]=[\phi(b),j+1]$.  They are strictly nested:
every $\phi(b)$ satisfies $p\phi(b)=\phi(b)$, where
$p=\phi(1_B)\ne1_B$, so $1_B$ cannot belong to $\phi(B)$.
Thus $[1_B,j+1]\notin M_j$, and
\begin{equation}\label{eq:strict-stages}
 X_0\subsetneq X_1\subsetneq X_2\subsetneq\cdots.
\end{equation}
The algebra $B$ is countable, being a quotient of the free associative
algebra over the finite field $k$ on finitely many generators.
Consequently $X=\bigcup_{j\geq0}X_j$ is countable; the strict inclusions
make it infinite.

\emph{Faithfulness of the affine part.}
Take a nonidentity element of $E_\infty$, represented by a matrix
$g\ne I$ at stage $j$.  Choose a nonzero entry $r$ of $g-I$, say in
row $a$ and column $b$.  At stage $j+1$, let $v$ have $b_*$ in
coordinate $b$ and zero in the other coordinates.  The $a$th coordinate
of $F(g)v-v$ is
\[
 \psi(r)\cdot b_*\ne0
\]
by Theorem~\ref{thm:algebra-package}.  It remains nonzero in the direct
limit because $\phi$ is injective.  This proves faithfulness of the
linear action.  If an affine element acts identically, evaluating it
at zero first makes its translation zero; linear faithfulness then
makes the element trivial.

\emph{Faithfulness of the whole group.}
Every affine limit element and its inverse have representatives at
some common stage.  At every later stage their translation vectors,
matrices, and inverse matrices act within that stage.  Hence, for each
affine limit element $a$,
\begin{equation}\label{eq:eventual-stage-preservation}
 a(X_j)=X_j\qquad\text{for all sufficiently large }j.
\end{equation}
In contrast,
\[
 t^k(X_j)=X_{j+k}
 \qquad\text{whenever }j\geq0\text{ and }j+k\geq0.
\]
For $k\ne0$, this differs from $X_j$ by
\eqref{eq:strict-stages}.  Thus no nonzero power of $t$ agrees, as a
permutation, with an affine limit element.  If $a t^k$ acts trivially,
then $t^k$ acts as $a^{-1}$, forcing $k=0$ by
\eqref{eq:eventual-stage-preservation}.  Affine faithfulness then gives
$a=1$.

\emph{Two-transitivity.}
We show that $E_\infty$ is transitive on the nonzero vectors of $X$.
Given two such vectors, choose a common representing stage and advance
once more.  All their coordinates now lie in $A$, since
$\phi(B)\subseteq A$, and the vectors remain nonzero.  We work in this
one stage, with the same simple algebra $A$ and its identity $e$.

If a vector has a nonzero coordinate $m_j\in A$, simplicity gives
$A m_j A=A$: the finite sums of elements $a m_j b$ form a nonzero
two-sided ideal.  Thus, for any $d\in A$, there are finitely many
$a_\nu,b_\nu\in A$ with
\[
 d=\sum_\nu a_\nu m_j b_\nu.
\]
The coefficient $r=\sum_\nu a_\nu\otimes b_\nu^{\op}\in R$ then
satisfies $r\cdot m_j=d$.  For $l\ne j$, an elementary matrix in row
$l$, column $j$, with coefficient $r$ adds $d$ to coordinate $l$ and
changes no other coordinate.  It can therefore set that coordinate to
any prescribed value in $A$.

If $j\ne6$, set the last coordinate to $e$.  If $j=6$, first set
coordinate $1$ to $e$, then use it to set the last coordinate to $e$.
Use the last coordinate to kill all the others.  This sends the vector
to $(0,0,0,0,0,e)$.  Applying the procedure to both original vectors
sends them to the same vector at the same stage, proving the asserted
linear transitivity.  Translations, followed by these linear maps and
another translation, take any ordered pair of distinct points to any
other such pair.  The action is two-transitive.

\emph{Point stabilizers.}
Both $E_\infty$ and $t$ fix zero, so the semidirect-product description
gives
\[
 \operatorname{Stab}_\Gamma(0)=E_\infty\rtimes\langle t\rangle
 \cong\Gamma/X.
\]
This is a finitely generated group, since $\Gamma$ is finitely
presented.  More explicitly, $t$ and the images of a finite generating
set of the base Steinberg group generate it: conjugating by positive
powers of $t$ reaches every linear stage.  All other point stabilizers
are conjugate to this one, since translations act transitively.
\end{proof}

\subsection{Embedding the input group and finishing the proof}

The last step is to put the original group into the elementary matrix
part of $\Gamma$.  Here the commutator property of
Lemma~\ref{lem:group-input} is essential.

\begin{proof}[Proof of Theorem~\ref{thm:envelope}]
Let $\widehat G$ be the auxiliary group of
Lemma~\ref{lem:group-input}.  Theorem~\ref{thm:algebra-package} supplies
an injective homomorphism $\iota:\widehat G\to A^\times$, where units
in $A$ have identity $e$.  Define
\[
 \nu(h)=1_B-e+\iota(h)\qquad(h\in\widehat G).
\]
The element $1_B-e$ annihilates $A$ on both sides.  Consequently
\[
 \nu(h)\nu(h')=\nu(hh'),\qquad
 \nu(h)^{-1}=\nu(h^{-1}).
\]
Injectivity of $\iota$ proves injectivity of $\nu$.  Composing with the
injective unital map $B\to R$, $b\mapsto b\otimes1_B$, gives a group
embedding $\bar\nu:\widehat G\to R^\times$.  The tensor map is
injective because $B$ is a nonzero vector space over $k$.

For $u\in R^\times$, let $D_{ij}(u)$ be the diagonal matrix with
entries $u,u^{-1}$ in positions $i,j$ and $1$ elsewhere.  Equation~
\eqref{eq:balanced-diagonal} shows that $D_{ij}(u)\in\E_6(R)$.
Let $s$ be the stable letter in $\widehat G$ and set
\[
 a_g=\bar\nu((g,g)),\qquad b=\bar\nu(s).
\]
By Lemma~\ref{lem:group-input},
$[a_g,b]=\bar\nu((g,1))$.  Taking the commutator of balanced diagonal
matrices on two different pairs of coordinates gives
\[
 [D_{12}(a_g),D_{13}(b)]
 =\operatorname{diag}\bigl(\bar\nu((g,1)),1,1,1,1,1\bigr).
\]
Thus the injective homomorphism given by the diagonal on the right
has image in $\E_6(R)$.  Since the transition $F$ is injective, this
embeds $G$ in $E_\infty$, and hence in $\Gamma$ by
Lemma~\ref{lem:affine-limit}.  Proposition~\ref{prop:two-transitive-action}
and the finite presentation of $\Gamma$ give all the remaining
assertions of Theorem~\ref{thm:envelope}.
\end{proof}

\begin{proof}[Proof of Theorem~\ref{thm:main}]
Suppose first that $G$ is finitely generated and has decidable word
problem.  Theorem~\ref{thm:envelope} embeds it in a finitely presented
group $\Gamma$ with a faithful two-transitive action on a countably
infinite set and finitely generated point stabilizers.  There is exactly one orbit on unordered
two-element subsets.  These are all the hypotheses of
Theorem~\ref{thm:zaremsky}, which embeds $\Gamma$ in a finitely
presented simple group $S$.  The group $S$ is nontrivial, since
$\Gamma$ acts transitively on an infinite set.  This proves the
required direction, including trivial and finite input groups.

For the converse, we give the classical decision argument attributed
by Boone and Higman to Kuznetsov \cite[p.~52]{BH74}.
Let $S=\langle Y\mid\mathcal R\rangle$ be a nontrivial
finitely presented simple group.  On an input word $w$, run the
following two searches in parallel:
\begin{enumerate}[label=(\roman*)]
 \item enumerate presentation proofs that $w=1$ follows from
       $\mathcal R$;
 \item enumerate, for every $y\in Y$, a presentation proof that
       $y=1$ follows from $\mathcal R\cup\{w\}$, stopping when
       a proof has been found for every generator.
\end{enumerate}
Consequences of a finite group presentation are recursively enumerable,
by enumerating finite products of conjugates of relators and their
inverses.  If $w=1$ in $S$, the first search succeeds, whereas the
second cannot succeed for every generator because $S$ is nontrivial.
If $w\ne1$, its normal closure in the simple group $S$ is all of $S$;
thus the second search succeeds for each of the finitely many
generators, whereas the first never succeeds.  Exactly one search
terminates, deciding the word problem in $S$.

For a finitely generated subgroup of $S$, fix words in $Y$ representing
its generators and substitute these words into each input.  This
transfers the decider to that subgroup.  Finally, the same substitution
argument between any two finite generating sets of a group proves
that decidability of its word problem is independent of the chosen
finite generating set.
\end{proof}
